\begin{workedexample}[Worked Example: Axial-mode helix gain]
An axial-mode helix radiates a circularly polarized beam along its axis
when the circumference $C \approx \lambda$. Kraus's gain estimate is
\[ G \approx 15\,N\,\frac{C^2}{\lambda^2}\,\frac{S}{\lambda}. \]
For $N = 10$ turns, $C = \lambda$, and spacing $S = 0.25\lambda$:
$G \approx 15\cdot10\cdot1\cdot0.25 = 37.5 = 15.7\ \text{dBi}$. More turns give
more gain and a narrower axial beam.
\end{workedexample}

\begin{keytakeaways}
\begin{itemize}[leftmargin=1.4em,itemsep=2pt]
\item Axial mode needs $C \approx \lambda$ and pitch $\sim$12--14$^{\circ}$; it is naturally circularly polarized.
\item Kraus: $G \approx 15N(C/\lambda)^2(S/\lambda)$; more turns $N$ raise gain.
\item Input impedance is near-resistive ($\sim140\,C/\lambda\ \Omega$), matched with a taper.
\end{itemize}
\end{keytakeaways}

\begin{exercises}
\item What circumference (in $\lambda$) puts a helix in axial mode?
\item Doubling the turns changes the Kraus gain by how many dB?
\item What polarization does an axial-mode helix radiate?
\end{exercises}

\furtherreading{Kraus~\cite{kraus} (ch.~7); Balanis~\cite{balanis}.}
